\documentclass[10pt,a4paper]{amsart}
\usepackage[margin=19mm]{geometry}
\usepackage{amsmath,amssymb,mathtools}
\usepackage[scr=rsfs,cal=euler]{mathalfa}
\usepackage{xfrac}
\usepackage[unicode,hidelinks,bookmarks=false]{hyperref}
\newcommand{\ie}{i.e.\ }
\newcommand{\cf}{cf.\ }
\newcommand{\resp}{respectively\ }
\newcommand{\Subsection}{\subsection}
\providecommand{\nobreakdash}{\nobreak\hbox{-}}
\setlength{\emergencystretch}{2em}
\multlinegap0pt
\usepackage{amssymb}
\usepackage{tikz}
\usepackage{float}
\usepackage{enumerate}
\newtheorem{lemma}{Lemma}
\title{On the difference of the intersection power graph and the power graph of a finite group}
\author{Sudip Bera, Peter J. Cameron}
\date{Source: arXiv:2509.03919v2 (CC BY 4.0)}
\begin{document}
\maketitle

\noindent\emph{Source and licence.} On the difference of the intersection power graph and the power graph of a finite group. Version arXiv:2509.03919v2 (CC BY 4.0), distributed by arXiv under the Creative Commons Attribution 4.0 International licence, http://creativecommons.org/licenses/by/4.0/, as recorded in the arXiv licence metadata for this exact version and shown on the abstract page https://arxiv.org/abs/2509.03919v2. Full licence notice: \texttt{licenses/arxiv-2509.03919-CC-BY-4.0.txt}.\par\smallskip

\noindent\textit{Editorial scope.} Counted unit: the article's lemma lemma:diam(B(G))-g-cyclic (source0.tex lines 437--442). The article's complete original proof of it is delivered with it (source0.tex lines 444--494, one \texttt{proof} environment of 51 lines). No mathematics is written, repaired or re-derived: the statement and the proof are the article's own lines, in the article's own notation, and no retained line is substituted or re-flowed.  No further source-local context is needed: every cross-reference of every retained line resolves inside the two delivered spans.  Nothing was omitted from any retained span, so the package carries no omission record.  No retained line contains a citation, so the wrapper carries no bibliography and no bibliography entry.  No retained line contains a figure environment, an image-inclusion command or drawing code, so no image is delivered and the package declares no asset dependency.  Editorial formatting only, in addition: an AMS \texttt{amsart} preamble that restates the article's own declarations and macros used by the retained lines, namely the environments \texttt{lemma} on separate counters, together with the standard packages those lines need.  The retained environments print in this document as Lemma 1 in order of appearance, whereas the article prints the same items with the numbering of its own full document.  The complete native source of the article as distributed in the arXiv e-print for this exact version is bundled unchanged as \texttt{sources/184450/source0.tex}.
\begin{lemma}\label{lemma:diam(B(G))-g-cyclic}
Let $G$ be a cyclic group of order $n$. Then the graph $\mathcal{B}(G) = \mathcal{B}(n)$ is empty if and only if $\pi(G) \leq 1$ or $G \cong \mathbb{Z}_{pq}$ for some distinct primes $p,q$. In all other cases, $\mathcal{B}(G)$ is connected with
\[
\operatorname{diam}\big(\mathcal{B}(G)\big) \leq 3.
\]
\end{lemma}

\begin{proof}
If $\pi(G) \leq 1$ or $G \cong \mathbb{Z}_{pq},$ then clearly $\mathcal{B}_1(G)$ is an empty graph.

Conversely,
let $x, y \mid n$ with $x, y \notin \pi(n) \cup \{1, n\}$ and $x \neq y$, where $\pi(n)$ is the set of prime divisors of a natural number $n$. We consider the following cases.

\medskip
\noindent\textbf{Case I:} $n = p_1 p_2 \cdots p_k$, $k \geq 3.$ 

\smallskip
\noindent\textit{Case IA.} If $\pi(x) \not\subseteq \pi(y)$ and $\pi(y) \not\subseteq \pi(x)$, then there exist primes $p_i, p_j$ with $p_i \mid x$, $p_i \nmid y$, $p_j \mid y$, $p_j \nmid x$. Hence
$x \sim p_ip_j \sim y.$

\smallskip
\noindent\textit{Case IB.} If $\pi(x) \subseteq \pi(y)$, then since $y \neq n$ we have $\pi(y) \neq \pi(n)$; choose $p_k \in \pi(n) \setminus \pi(y)$ and $p_i \in \pi(x)$. Then $x \sim p_ip_k \sim y.$

\medskip
\noindent\textbf{Case II:} $n = p_1^{\alpha_1}p_2^{\alpha_2}\cdots p_k^{\alpha_k}$, $k \geq 2$, with some $\alpha_i > 1$.

\smallskip
\noindent\textit{Case IIA.} Analogous to Case IA.

\smallskip
\noindent\textit{Case IIB.} If $\pi(x) \subseteq \pi(y) \subsetneq \pi(n)$, the argument is analogous to Case IB.

\smallskip
\noindent\textit{Case IIC.} Suppose $\pi(x) \subsetneq \pi(y) = \pi(n)$. Since $y \neq n$, there exists $k$ with $\alpha_k > 1$ and $p_k^{\alpha_k} \nmid y$; in particular $y \nmid x$.
\begin{itemize}
\item If $x \nmid y$, then $x \sim y$ directly.
\item If $x \mid y$ and $x = p_k^{r_k}$ with $1 < r_k < \alpha_k$, then for any $i \neq k$,
\[
x = p_k^{r_k} \sim p_ip_k \sim p_k^{\alpha_k} \sim y.
\]
\item If $x \mid y$ and there exists $p_i \mid x$ with $i \neq k$, then
\[
x \sim p_ip_k \sim p_k^{\alpha_k} \sim y \quad \text{if } p_k \nmid x, \qquad\text{and}\qquad x \sim p_k^{\alpha_k} \sim y \quad \text{if } p_k \mid x.
\]
\end{itemize}

\smallskip
\noindent\textit{Case IID.} Suppose $\pi(x) = \pi(y) = \pi(n)$.
\begin{itemize}
\item If $x \nmid y$ and $y \nmid x$, then $x \sim y$ directly.
\item If $x \mid y$ and $y \nmid x$: since $y \neq n$, there exists $k$ with $\alpha_k > 1$ and $p_k^{\alpha_k} \nmid y$, whence also $p_k^{\alpha_k} \nmid x$. Then
\[
x \sim p_k^{\alpha_k} \sim y.
\]
\end{itemize}

In every case, $d(x,y) \leq 3$, which proves the bound.
\end{proof}

\end{document}
